\lyseis{2.1}{Η έννοια του κλάσματος}
\sol{1} α) $\frac{3}{8}$ \quad β) $\frac{5}{6}$ \quad γ) $\frac{1}{4}$ \quad δ) $\frac{7}{10}$.
\sol{2} Χωρίζεις κάθε ορθογώνιο σε ίσα μέρη, με το ίδιο πλάτος: σε 5 για το α) και χρωματίζεις 2, σε 4 για το β) και χρωματίζεις 3, σε 6 για το γ) και χρωματίζεις 5.
\sol{3} α) $\frac{3}{7}$ \quad β) $\frac{5}{4}$ \quad γ) $\frac{1}{10}$ \quad δ) $\frac{7}{9}$. Διαβάζονται: «τέσσερα ένατα», «ένα δωδέκατο», «έντεκα έκτα», «τρία εκατοστά».
\sol{4} $5:8=\frac{5}{8}$, $3:11=\frac{3}{11}$, $10:3=\frac{10}{3}$, $7:7=\frac{7}{7}$ (που ισούται με 1). Και $\frac{9}{4}=9:4$, $\frac{1}{6}=1:6$, $\frac{20}{5}=20:5$ (που ισούται με 4).
\sol{5} $7=\frac{7}{1}=\frac{28}{4}=\frac{70}{10}$ (γιατί $7\cdot 4=28$ και $7\cdot 10=70$). Και $0=\frac{0}{5}$.
\sol{6} Μικρότερα από το 1: $\frac{3}{4}$, $\frac{5}{12}$, $\frac{1}{100}$. Ίσα με το 1: $\frac{7}{7}$, $\frac{20}{20}$. Μεγαλύτερα από το 1: $\frac{10}{9}$, $\frac{18}{17}$.
\sol{7} α) $\frac{25}{60}$ \quad β) $\frac{250}{1000}$ \quad γ) $\frac{2}{7}$ \quad δ) $\frac{3}{12}$.
\sol{8} α) $20:4\cdot 3=15$ \quad β) $35:5\cdot 2=14$ \quad γ) $48:6\cdot 5=40$ \quad δ) $200:10\cdot 7=140$.
\sol{9} α) $14:2\cdot 3=21$ \quad β) $27:3\cdot 5=45$ \quad γ) $36:4\cdot 9=81$ \quad δ) $56:7\cdot 8=64$.
\sol{10} Το ένα πέμπτο των 20 είναι 4, άρα τα τρία πέμπτα είναι $3\cdot 4=12$ βολές. Έχασε $20-12=8$ βολές, δηλαδή τα $\frac{2}{5}$ ($1-\frac{3}{5}$).
\sol{11} Το ένα ένατο των 36 είναι 4, άρα τα $\frac{5}{9}$ είναι $5\cdot 4=20$ χρωματιστά. Τα υπόλοιπα είναι $36-20=16$, δηλαδή τα $\frac{4}{9}$ των μολυβιών.
\sol{12} Το ένα όγδοο των 2.000 γρ. είναι 250 γρ., άρα τα $\frac{3}{8}$ είναι $3\cdot 250=750$ γραμμάρια.
\sol{13} Ο Βασίλης: $27:9\cdot 4=12$ ψήφους. Η Ζωή: $27:3=9$ ψήφους. Η Μυρτώ: $27-12-9=6$ ψήφους, δηλαδή τα $\frac{6}{27}$, που είναι ίσα με τα $\frac{2}{9}$ ($\frac{9}{9}-\frac{4}{9}-\frac{3}{9}$).
\sol{14} Τα $\frac{5}{8}$ είναι 150 χλμ., άρα το ένα όγδοο είναι $150:5=30$ χλμ. και όλη η διαδρομή $8\cdot 30=240$ χλμ. Μένουν $240-150=90$ χλμ.
\sol{15} Τα $\frac{3}{7}$ είναι 18 λίτρα, άρα το ένα έβδομο είναι $18:3=6$ λίτρα και ολόκληρο το δοχείο χωράει $7\cdot 6=42$ λίτρα.
\sol{16} $\frac{15}{60}$, $\frac{20}{60}$ και $\frac{45}{60}$ της ώρας. Η σχολική μέρα είναι τα $\frac{6}{24}$ του εικοσιτετραώρου.
\sol{17} Αν κόψεις και τις δύο πίτες σε 9 κομμάτια, το $\frac{1}{9}$ κάθε πίτας είναι ένα κομμάτι. Παίρνεις ένα από κάθε πίτα, άρα 2 κομμάτια των $\frac{1}{9}$, δηλαδή $\frac{2}{9}$ μιας πίτας. Με υπολογισμό: $2\cdot\frac{1}{9}=\frac{2}{9}$.
\sol{18} Το συμπέρασμα είναι σωστό: $\frac{3}{5}<\frac{3}{4}$. Αν κόψεις μια πίτα σε 4 κομμάτια, κάθε κομμάτι είναι μεγαλύτερο από ό,τι αν την κόψεις σε 5 (περισσότερα κομμάτια σημαίνει μικρότερα κομμάτια). Τρία μεγάλα κομμάτια ($\frac{3}{4}$) είναι περισσότερο φαγητό από τρία μικρά ($\frac{3}{5}$).
\sol{19} Για να είναι το $\frac{7}{\nu}>1$ πρέπει $\nu<7$. Για να είναι $<2$ πρέπει $7<2\nu$, δηλαδή $\nu\geq 4$. Άρα $\nu=4,\,5,\,6$ και τα κλάσματα είναι τα $\frac{7}{4}$, $\frac{7}{5}$, $\frac{7}{6}$: τρία.
\sol{20} Αν $\frac{5}{0}=5:0=x$, τότε $5=0\cdot x$, που είναι αδύνατο, γιατί το $0\cdot x$ είναι πάντα 0. Δεν υπάρχει αριθμός που να είναι το $5:0$. Γι' αυτό ο παρονομαστής δεν μπορεί να είναι 0.
\sol{21} Χωρίζω ολόκληρη την πίτα σε 6 ίσα κομμάτια. Η μισή πίτα είναι 3 κομμάτια. Τα $\frac{2}{3}$ της μισής είναι 2 από αυτά τα 3 κομμάτια, δηλαδή $\frac{2}{6}$ της ολόκληρης πίτας. (Το $\frac{2}{6}$ είναι ίσο με το $\frac{1}{3}$, όπως θα δούμε στην παράγραφο 2.2.)
\sol{22} $\mathrm{ΓΔ}=12:4\cdot 3=9$ εκ. και $\mathrm{ΕΖ}=12:6\cdot 5=10$ εκ.
\sol{23} Ενδεικτικά: μια σοκολάτα με 12 τετράγωνα, από τα οποία έφαγες 3: $\frac{3}{12}$. Ένα ημερολόγιο: ο Μάρτιος είναι $\frac{1}{12}$ του έτους. Μια αυγοθήκη των 6 με 2 αυγά: $\frac{2}{6}$.
